
\subsubsection{Research Questions}

The experiments address three questions:

\begin{enumerate}
\item Does CoT prompting reliably produce reasoning chains with epistemic hedging markers?
\item Are hedging markers structurally positioned to influence confidence generation?
\item Does hedging marker presence correlate with lower verbalized confidence?
\end{enumerate}

\subsubsection{Evaluation Metrics}

\begin{itemize}
\item \textbf{Hedging Marker Count:} Integer count of 17 predefined markers per output.
\item \textbf{Hedging Presence Rate:} Proportion of outputs containing at least one marker.
\item \textbf{Verbalized Confidence:} Numerical confidence (0-100\%) extracted from output.
\item \textbf{Reasoning Rate:} Proportion of outputs containing multi-step reasoning.
\item \textbf{Step Count:} Number of reasoning steps identified via sentence segmentation.
\item \textbf{Spearman Correlation:} Rank correlation between hedging count and confidence.
\end{itemize}

\subsubsection{Prompting Conditions}

The primary evaluation used the CoT+Confidence condition (CoT reasoning + answer + confidence). A baseline condition (direct answer only) was included for reasoning rate verification. The planned token-padding control condition (random filler tokens + answer + confidence) was not executed due to resource constraints.
